% !TeX root = ../thuthesis-example.tex

\chapter{绪论}

\section{研究背景与意义}

推荐系统是现代互联网平台的核心技术之一，旨在从海量候选物品中为用户筛选出可能感兴趣的内容。随着电子商务、社交媒体和内容平台的快速发展，推荐系统在提升用户体验和商业价值方面发挥着日益重要的作用。

序列推荐（Sequential Recommendation）是推荐系统领域的重要研究方向，其核心思想是根据用户的历史交互序列预测其下一个可能交互的物品。传统的序列推荐方法，如基于循环神经网络的GRU4Rec\cite{hidasi2016session}和基于自注意力机制的SASRec\cite{kang2018self}，通过学习物品ID的嵌入表示来建模用户偏好的动态变化。然而，这类方法依赖于独立的物品ID嵌入，无法建模物品之间的语义关联，且对新物品的泛化能力有限。

近年来，生成式推荐（Generative Recommendation）作为一种新范式受到了广泛关注\cite{rajput2023recommender,zheng2024adapting}。其核心思想是：首先为每个物品分配一个由多层离散编码组成的语义ID（Semantic ID），然后将推荐问题转化为序列到序列的生成任务——给定用户历史交互物品的语义ID序列，自回归地生成下一个物品的语义ID。与传统方法相比，生成式推荐具有以下优势：（1）语义ID编码了物品的内容特征，语义相似的物品拥有相近的编码，有助于知识迁移；（2）生成式框架天然支持开放词汇，对新物品具有更好的泛化性；（3）多层级的编码结构实现了从粗到细的物品检索。

在生成式推荐框架中，语义ID的构建质量直接决定了推荐性能的上限。目前主流的语义ID构建方法基于残差量化变分自编码器（RQ-VAE）\cite{lee2022autoregressive,rajput2023recommender}，将物品的文本嵌入通过编码器映射到低维隐空间，再经过多层残差向量量化生成离散编码。然而，现有方法在模态利用和信号融合方面存在明显不足：

第一，\textbf{模态利用单一}。主流方法仅利用文本语义信息构建语义ID。以TIGER\cite{rajput2023recommender}为代表的方法使用预训练语言模型提取物品文本嵌入，然后直接进行向量量化。然而，在电商等场景中，物品的视觉外观往往是影响用户决策的关键因素，忽视图像信息会导致语义ID无法全面描述物品特征。

第二，\textbf{协同过滤信号融合不够直接}。LETTER\cite{zheng2024adapting}率先尝试在语义ID构建中引入协同过滤（Collaborative Filtering, CF）信号，通过在RQ-VAE的训练损失中添加InfoNCE对比损失来对齐量化后的隐向量与CF嵌入。这种通过辅助损失间接融合的方式需要额外引入对比损失函数并仔细调节其权重，融合程度难以控制等问题。

第三，\textbf{缺乏有效的多模态融合机制}。现有的语义ID构建工作尚未探索如何将文本、图像和协同过滤三种异构模态有效融合到统一的语义ID表示中。简单的特征拼接或加权平均无法捕捉模态间的复杂交互关系。

\section{研究内容与创新点}

针对上述问题，本文提出了COMET（Collaborative-guided Multimodal Encoding Tokenizer），一种基于协同过滤引导的多模态语义ID构建方法。本文的主要研究内容和创新点包括：

（1）\textbf{提出了CF-as-Query跨注意力融合机制}。COMET将协同过滤嵌入作为注意力查询（Query），将文本多token投影和图像嵌入构成的键值序列作为键值对（Key-Value），通过多层堆叠的跨注意力和前馈网络实现三模态深度融合。文本嵌入被分割为4个子空间分别投影，形成4个KV token，使注意力能够细粒度地选择文本信息。这一设计的直觉在于：协同过滤信号反映了用户群体的行为偏好，用其作为查询可以自适应地从文本和视觉特征中提取对推荐最有价值的信息。与LETTER的辅助损失方式相比，COMET将协同过滤信号通过模型架构直接融入编码过程，无需引入额外的对比损失函数及其权重调节。

（2）\textbf{提出多目标重建策略促进多模态信息均衡编码}。COMET采用三个独立的解码器分支，以加权方式同时重建文本、图像和协同过滤三种模态的嵌入。文本嵌入作为维度最高、信息量最丰富的模态赋予最大权重，图像和协同过滤嵌入的重建起正则化作用，迫使隐向量在编码文本语义的同时保留视觉和协同过滤信号，实现多模态信息在隐空间中的均衡编码。

（3）\textbf{构建了完整的多模态生成式推荐实验框架}。本文搭建了从原始数据处理、多模态嵌入提取、语义ID构建到生成式推荐模型训练和评估的完整实验流水线，实现了RQ-KMeans、RQ-VAE、LETTER和COMET四种语义ID构建方法的公平对比，并在TIGER和OneRec两种代表性生成式推荐模型上验证了COMET的有效性。

\section{论文组织结构}

本文的组织结构如下：

第一章为绪论，介绍研究背景、问题提出和主要创新点。

第二章为相关工作，综述序列推荐、生成式推荐、语义ID构建和多模态推荐的相关研究。

第三章为方法，详细描述COMET的整体框架，包括多模态嵌入提取、CF-as-Query跨注意力融合模块、多目标重建、残差向量量化、训练策略和碰撞优化等。

第四章为实验，在Amazon Beauty数据集上进行全面的实验评估，包括主实验、消融实验和语义ID质量分析。

第五章为总结与展望，总结本文的工作并讨论未来的研究方向。
